La génération des fichiers météorologiques, également connus sous le nom de fichiers "bulk". Ces fichiers sont essentiels pour améliorer la précision du modèle CROCO. En effet, ils fournissent des informations détaillées sur diverses composantes atmosphériques, telles que la radiation, la température, la vitesse du vent, et bien d'autres. Ces composantes jouent un rôle crucial dans la modélisation des interactions océan-atmosphère, et leur intégration permet d'obtenir des simulations plus réalistes et précises. Nous détaillerons ce processus dans le sous-chapitre suivant.
